% ===========================================================================
%  sviluppatore/06-dashboard/03-job.tex
% ===========================================================================
\section{Il lavoro lungo: job e registro in diretta}\label{sez:job}

Un backtest dura minuti o ore. Una richiesta HTTP non può aspettarlo, quindi
tutto ciò che è lungo passa dal \simbolo{JobManager} di
\file{execution/job\_manager.py}: la rotta restituisce subito un
identificativo, il lavoro continua in un thread, e il registro arriva al
browser man mano che viene prodotto.

\begin{center}
\begin{tabular}{@{}L{40mm}L{98mm}@{}}
\toprule
\textbf{Rotta} & \textbf{Che cosa fa} \\
\midrule
\rigaTenue
\rotta{GET}{/api/jobs} & Elenco. \\
\rotta{GET}{/api/job/<id>} & Stato di uno. \\
\rigaTenue
\rotta{GET}{/api/job/<id>/stream} & \textbf{SSE}: le righe di registro, in
  diretta. \\
\rotta{POST}{/api/job/<id>/cancel} & Annulla. \\
\rigaTenue
\rotta{DELETE}{/api/job/<id>} & Toglie dall'elenco un job concluso. \\
\rotta{DELETE}{/api/jobs/conclusi} & Svuota tutti i conclusi. \\
\bottomrule
\end{tabular}
\end{center}

\subsection{Gli stati}

\cod{PENDING} \freccia{} \cod{RUNNING} \freccia{} \cod{SUCCESS},
\cod{FAILED} o \cod{CANCELLED}. Sono un \cod{str, Enum}, quindi il valore
serializzato in JSON è direttamente la stringa minuscola.

\subsection{Il registro in diretta}

\begin{codiceInterno}{python}{Lo stream SSE}{dashboard/app.py:3720}
    def api_job_stream(job_id):
        """Server-Sent Events stream of log lines."""
        job = job_manager.get(job_id)
        if job is None:
            return jsonify({"error": "Job not found"}), 404

        def event_stream():
            # Manda subito le log lines gia' esistenti
            existing = list(job.log_lines)
            for line in existing:
                yield f"data: {line}\n\n"

            # Subscribe via queue per nuove lines
            q = Queue(maxsize=10000)
            def listener(line):
                try:
                    q.put_nowait(line)
                except Exception:
                    pass
            unsub = job.subscribe(listener)
\end{codiceInterno}

\noindent
Due dettagli meritano di essere notati, perché sono la ragione per cui
questo codice regge:

\begin{description}[leftmargin = 1.6em, style = nextline, itemsep = 4pt]
  \item[Prima lo storico, poi la sottoscrizione]
    Chi apre la pagina a metà di un backtest vede subito tutto ciò che è già
    successo, e poi il resto. Sottoscrivere senza mandare lo storico
    mostrerebbe una finestra vuota fino alla riga successiva.
  \item[La coda è limitata e il \cod{put} non blocca]
    \cod{Queue(maxsize=10000)} con \cod{put\_nowait} dentro un
    \cod{try/except} che ignora il fallimento. Se un browser smette di
    leggere, la sua coda si riempie e le righe nuove vengono scartate --- ma
    il backtest \emph{non} rallenta. L'alternativa (bloccare sulla coda
    piena) farebbe sì che una scheda lasciata aperta fermi un calcolo.
\end{description}

\begin{nota}{Il registro è limitato a \num{5000} righe}
\cod{deque(maxlen=5000)}: un backtest lungo ne produce molte di più, e le
prime scorrono via. Per conservarle tutte si guarda il file di registro, non
l'interfaccia.
\end{nota}

\subsection{Un lavoro per volta}

Il \simbolo{JobManager} nasce con \cod{max\_workers=1}. È il default
dichiarato del progetto, e vale sia per l'API FastAPI sia per la dashboard.
Due backtest in parallelo si contenderebbero la stessa cache dei dati e gli
stessi file di uscita.

\begin{attenzione}{Alzare \cod{max\_workers} non è una ottimizzazione}
La serializzazione non è un limite di prestazioni da togliere: è ciò che
rende i risultati riproducibili. Due job concorrenti che scrivono sotto
\file{output/strategies/} possono produrre uno stream mezzo di un
calcolo e mezzo dell'altro, e il risultato non dà errore --- è un
\file{parquet} valido con dentro numeri di due esecuzioni diverse.
\end{attenzione}

\subsection{Le notifiche}

\rotta{GET}{/api/notifications} alimenta la campanella, e mette insieme due
sorgenti: i job reali del \simbolo{JobManager} e i messaggi che arrivano dal
sito. È il motivo per cui la campanella mostra sia «backtest finito» sia
«hai una risposta su una contestazione»: una rotta, due origini.
